\section{10BASE-T1S / PLCA Validation Plan}
\label{sec:t1s-plca}

T1S validation is a first-class deliverable. ``Packets were successfully transmitted'' is not sufficient validation anywhere in this program -- the network itself must be characterized, not just exercised.

\subsection{What must be characterized}

\begin{enumerate}[itemsep=0.05em]
\item Physical link bring-up
\item PLCA operation
\item PLCA coordinator behavior
\item Node-ID configuration
\item Transmit opportunities
\item End-to-end latency
\item Node-to-node latency
\item Jitter
\item Packet loss
\item Bus utilization
\item Maximum sustained clean traffic rate
\item Cold-start time
\item Node disconnect behavior
\item Node reconnect time
\item Power-cycle recovery
\item Termination sensitivity
\item Cable/topology sensitivity
\item Error recovery
\item Node power consumption
\item Maximum node count actually tested
\end{enumerate}

If practical, additionally characterize gPTP/time synchronization and endpoint synchronization accuracy -- LAN8660/LAN8661 support optional gPTP per Microchip's public documentation (Section~\ref{sec:hardware-roles}), so this is a plausible stretch measurement, not a required one.

\begin{decision}[title={Datasheet numbers are not measured results}]
Every item above is reported from an actual bench measurement on Alfred's own network, with the procedure and equipment stated (Section~\ref{sec:measurement}). A datasheet's typical/maximum figure is a design input, never a substitute for a result.
\end{decision}

\subsection{PLCA-specific proof}

Use an oscilloscope or network capture to prove actual PLCA operation, not just infer it from successful application-level traffic:

\begin{diagram}
    coordinator beacon
        |
    node transmit opportunity 1
        |
    node transmit opportunity 2
        |
    node transmit opportunity 3
\end{diagram}

Test progression, each stage building on the last:

\begin{itemize}[itemsep=0.1em]
\item Two-node T1S (minimum viable PLCA proof).
\item Three-node T1S.
\item Clicker/LAN8651, LAN8660, and LAN8661 together on the \emph{same} segment -- this is the configuration the mandatory Greenfield proof (Section~\ref{sec:arch-a-proof}) actually requires, and it is validated at the network-characterization level here, not assumed from the two/three-generic-node results above.
\end{itemize}

The mandatory November Greenfield network must demonstrate PLCA with both LAN8660 and LAN8661 active simultaneously -- a passing two-node COTS-only PLCA test (Section~\ref{sec:cots-bench}) establishes that the network layer works in general; it does not by itself establish that LAN8660/LAN8661 behave correctly as PLCA followers, which is a separate, required test.
